\documentclass[10pt,a4paper]{amsart}
\usepackage[margin=19mm]{geometry}
\usepackage{amsmath,amssymb,mathtools}
\usepackage[scr=rsfs,cal=euler]{mathalfa}
\usepackage{xfrac}
\usepackage[unicode,hidelinks,bookmarks=false]{hyperref}
\newcommand{\ie}{i.e.\ }
\newcommand{\cf}{cf.\ }
\newcommand{\resp}{respectively\ }
\newcommand{\Subsection}{\subsection}
\providecommand{\nobreakdash}{\nobreak\hbox{-}}
\setlength{\emergencystretch}{2em}
\multlinegap0pt
\usepackage{amssymb}
\usepackage{tikz}
\newtheorem{lemma}{Lemma}
\title{Anti‑Ramsey Numbers for Spanning Linear Forests of 3‑Vertex Paths and Matchings}
\author{Ali Ghalavand, Xueliang Li}
\date{Source: arXiv:2509.25949v3 (CC BY 4.0)}
\begin{document}
\maketitle

\noindent\emph{Source and licence.} Anti‑Ramsey Numbers for Spanning Linear Forests of 3‑Vertex Paths and Matchings. Version arXiv:2509.25949v3 (CC BY 4.0), distributed by arXiv under the Creative Commons Attribution 4.0 International licence, http://creativecommons.org/licenses/by/4.0/, as recorded in the arXiv licence metadata for this exact version and shown on the abstract page https://arxiv.org/abs/2509.25949v3. Full licence notice: \texttt{licenses/arxiv-2509.25949-CC-BY-4.0.txt}.\par\smallskip

\noindent\textit{Editorial scope.} Counted unit: the article's lemma lm1 (source0.tex lines 162--172). The article's complete original proof of it is delivered with it (source0.tex lines 173--220, one \texttt{proof} environment of 48 lines). No mathematics is written, repaired or re-derived: the statement and the proof are the article's own lines, in the article's own notation, and no retained line is substituted or re-flowed.  No further source-local context is needed: every cross-reference of every retained line resolves inside the two delivered spans.  Nothing was omitted from any retained span, so the package carries no omission record.  No retained line contains a citation, so the wrapper carries no bibliography and no bibliography entry.  No retained line contains a figure environment, an image-inclusion command or drawing code, so no image is delivered and the package declares no asset dependency.  Editorial formatting only, in addition: an AMS \texttt{amsart} preamble that restates the article's own declarations and macros used by the retained lines, namely the environments \texttt{lemma} on separate counters, together with the standard packages those lines need.  The retained environments print in this document as Lemma 1 in order of appearance, whereas the article prints the same items with the numbering of its own full document.  The complete native source of the article as distributed in the arXiv e-print for this exact version is bundled unchanged as \texttt{sources/186043/source0.tex}.
\begin{lemma}\label{lm1}
For three positive integers \( k \), \( t \), and \( n \), where $k\geq1$, \(t \geq 2\), and \(n = 3k + 2t\), let \( c \) be an edge-coloring of \( K_n \) such that 
\[
|c(K_n)| \geq \frac{1}{2} (3k + 2t - 3)(3k + 2t - 4) + 2.
\]
If \( K_{n-3} \) is a subgraph of \( K_n \) such that 
$|c(K_{n-3})|$ is maximized, then it holds that
\[
|c(K_{n-3})| \geq \frac{1}{2}(3k + 2t - 6)(3k + 2t - 7) + 2.
\]
\end{lemma}

\begin{proof}
  For three positive integers \(k\), \(t\), and \(n\) where $k\geq1$, \(t \geq 2\), and \(n = 3k + 2t\), let \(c\) be an edge-coloring of \(K_n\) such that 
\[
|c(K_n)| \geq \frac{1}{2}(3k + 2t - 3)(3k + 2t - 4) + 2.
\]
Assume that \(G\) is a rainbow spanning subgraph of size \(|c(K_n)|\) of \(K_n\), and let \(K_{n-3}\) be a subgraph of \(K_n\) such that \(|E(K_{n-3}) \cap E(G)|\) is maximized. To prove our lemma, it suffices to demonstrate that 
\[
|E(K_{n-3}) \cap E(G)| \geq \frac{1}{2}(3k + 2t - 6)(3k + 2t - 7) + 2.
\]
To establish this,  we will analyze the situation in two distinct cases as follows:

\noindent
{\bf Case I:} For every vertex \(v\) of \(K_{n-3}\), it holds that
   \[
   |\{uv : u \in V(K_{n-3})\} \cap E(G)| \geq 3k + 2t - 6.
   \]
   In this case, since $k\geq1$ and \(t \geq 2\), it follows that 
   \[
   |E(K_{n-3}) \cap E(G)| \geq \frac{1}{2}(3k + 2t - 3)(3k + 2t - 6) \geq \frac{1}{2}(3k + 2t - 6)(3k + 2t - 7) + 2,
   \]
   which is as desired.
   
   \noindent
 {\bf Case II:} For a vertex \(v\) of \(K_{n-3}\), it holds that
   \[
   |\{uv : u \in V(K_{n-3})\} \cap E(G)| \leq 3k + 2t - 7,
   \]
   and it also holds that 
   \[
   |E(K_{n-3}) \cap E(G)| \leq \frac{1}{2}(3k + 2t - 6)(3k + 2t - 7) + 1.
   \]
   In this case, since \(|E(G)| = |c(K_n)| \geq \frac{1}{2}(3k + 2t - 3)(3k + 2t - 4) + 2\), it follows that 
   \[
   |E(G) - E(K_{n-3})| \geq 9k + 6t - 14.
   \]
   Thus, there exists a vertex \(x\) in \(V(K_n) - V(K_{n-3})\) such that 
   \[
   |\{ux : u \in V(K_{n-3})\} \cap E(G)| \geq 3k + 2t - 5.
   \]
   However, this leads to a contradiction regarding the maximality, because using the vertices \(v\) and \(x\), we can construct a graph \(K_{n-3}\) such that 
   \(|E(K_{n-3}) \cap E(G)|\) is larger.
   
Now, based on these two cases, we can conclude that 
\[
|E(K_{n-3}) \cap E(G)| \geq \frac{1}{2}(3k + 2t - 6)(3k + 2t - 7) + 2,
\]
therefore completing the proof.
\end{proof}

\end{document}
